%% ma: L12-B14-0104
%% lop: 12
%% chuong: 5
%% bai: 14
%% dang: 12.14.31
%% loai: TLN
%% muc: VD
%% dap_an: "36"
%% nguon: "Ảnh giáo viên gửi 07/11/2026 (Tư Duy Mở - tuduymo.com, đề tự luyện Phương trình mặt phẳng, đề lập phương trình mặt phẳng), Câu 22"
%% kien_thuc: "Bài 14 (phương trình mặt phẳng tổng quát, các trường hợp riêng), thể tích lăng trụ (lớp 11)"
%% trang_thai: chinh-thuc
%% ghi_chu: "Tài liệu gốc không kèm đáp án; đáp án do Claude giải và kiểm bằng sympy. Dạng đã duyệt 07/11/2026."
\begin{ex}
Trong không gian $Oxyz$, cho hai mặt phẳng lần lượt là $(\alpha)\colon 2x+z-6=0$ và $(\beta)\colon mx+(n-1)y+nz+4-n=0$, với $m,n\in\mathbb R$. Biết rằng hình đa diện giới hạn bởi $(\alpha)$, $(\beta)$ và $3$ mặt phẳng tọa độ là một lăng trụ đứng $(H)$. Hãy xác định thể tích của $(H)$.
\shortans{36}
\loigiai{$(\alpha)$ song song với $Oy$ và cùng với $(Oxy)$, $(Oyz)$ tạo thành mặt bên của một lăng trụ có cạnh bên song song $Oy$; hai đáy là $(Ozx)\colon y=0$ và $(\beta)$, nên $(\beta)\parallel(Ozx)$: $m=n=0$, $(\beta)\colon y=4$.
Đáy là tam giác vuông trong $(Ozx)$ với các cạnh góc vuông $3$ (trên $Ox$) và $6$ (trên $Oz$), diện tích $9$; chiều cao $4$. $V=9\cdot4=36$.}
\end{ex}
